\begin{frame}{Supervised Fine-Tuning (SFT)}
  \begin{center}
    \Large
    Supervised Fine-Tuning (SFT)
  \end{center}
\end{frame}


\begin{frame}{What is Supervised Fine-Tuning (SFT)?}
  \centering
  \includegraphics[width=0.60\paperwidth,height=0.60\paperheight,keepaspectratio]{assets/Lecture-8_Fine_Tuning_RLHF/slide_015.png}
\end{frame}

\begin{frame}{Supervised Fine-Tuning For Gemini LLM}
  \centering
  \includegraphics[width=0.60\paperwidth,height=0.60\paperheight,keepaspectratio]{assets/Lecture-8_Fine_Tuning_RLHF/slide_016.png}
\end{frame}

\begin{frame}{Datasets for Supervised Fine-Tuning (SFT)}
  \begin{itemize}
    \item \textbf{Instructional datasets:}
    \begin{itemize}
      \item \textbf{OpenAI:} InstructGPT (Anthropic Helpfulness data)
      \item \textbf{Stanford Alpaca:} 52k GPT-3 generated instructions
      \item \textbf{Others:} Dolly, ShareGPT, OASST, UltraChat
    \end{itemize}
    \item[\textcolor{red}{$\blacktriangleright$}] \textbf{Note:} Quality of supervision greatly affects model behavior.
  \end{itemize}
\end{frame}

\begin{frame}{Limitations of Supervised Fine-Tuning (SFT)}
  \begin{itemize}
    \item[\textcolor{red}{$\blacktriangleright$}] Cannot capture nuanced human preferences or values.
    \item[\textcolor{red}{$\blacktriangleright$}] May reinforce existing biases or hallucinations present in the data.
    \item[\textcolor{red}{$\blacktriangleright$}] Risk of overfitting, especially on synthetic or noisy datasets.
    \item[\textcolor{red}{$\blacktriangleright$}] Often leads to safe but bland and generic responses.
  \end{itemize}
\end{frame}
